% Figure: the five fifth-roots of unity at the corners of a regular pentagon
% inscribed in the unit circle, with the vertex at 2 pi / 5 highlighted as the
% source of the exact value cos(2 pi / 5) = (sqrt(5) - 1) / 4.
\begin{figure}[htbp]
\centering
\begin{tikzpicture}[scale=2.2, line cap=round, >=latex, font=\footnotesize]
  % unit circle and axes
  \draw[enggray, thin] (0,0) circle (1);
  \draw[->, engblack, thin] (-1.25,0) -- (1.3,0) node[right] {$\Re$};
  \draw[->, engblack, thin] (0,-1.25) -- (0,1.3) node[above] {$\Im$};
  % five roots at angles 2 pi k / 5, k = 0..4
  \foreach \k in {0,1,2,3,4} {
    \pgfmathsetmacro{\ang}{72*\k}
    \filldraw[engblue] (\ang:1) circle (0.028);
  }
  % pentagon edges
  \draw[engblue, thick] (0:1) -- (72:1) -- (144:1) -- (216:1) -- (288:1) -- cycle;
  % highlight k=1 vertex and its real projection
  \filldraw[engred] (72:1) circle (0.035);
  \node[engred, anchor=south west] at (72:1.03) {$e^{i 2\pi/5}$};
  \draw[engred, dashed] (72:1) -- ({cos(72)}, 0);
  \filldraw[engred] ({cos(72)}, 0) circle (0.02);
  \node[engred, anchor=north] at ({cos(72)}, -0.02) {$\cos\tfrac{2\pi}{5}$};
  % label k=0 root
  \node[engblue, anchor=west] at (0:1.03) {$1$};
  % central angle arc
  \draw[engamber] (0.32,0) arc (0:72:0.32);
  \node[engamber] at (36:0.46) {$\tfrac{2\pi}{5}$};
\end{tikzpicture}
\caption[The five fifth-roots of unity]{The five fifth-roots of unity, the
solutions of $z^{5} = 1$, at the corners of a regular pentagon inscribed in
the unit circle. Their common magnitude is one and their arguments are
$2\pi k/5$ for $k = 0,\dots,4$, spaced by $72^{\circ}$. The five roots sum to
zero, the same cancellation that balances a five-phase current set. The real
projection of the $k = 1$ vertex is the exact value
$\cos(2\pi/5) = (\sqrt{5} - 1)/4$, the algebraic number that the classical
straightedge-and-compass construction of the pentagon depends on. Reading the
constructible cosine off the root of unity is the route by which the roots of
unity connect to classical geometry.}
\label{fig:vol02ch03:fifth-roots}
\end{figure}
